\section{Reading the first baseline results}

First real walk-forward backtest (2026-07-14). Test: 12 months of PL load,
8{,}666 hours, day-ahead under the 09:00 D-1 cutoff. Weather input: ERA5
actuals --- optimistic; honest rerun with archived forecasts follows.

\subsection{The table}

\begin{center}
\begin{tabular}{lrrr}
\toprule
Model & MAPE (\%) & MAE (MW) & Skill vs naive \\
\midrule
ridge          & 4.00 & 713  & 0.31 \\
LASSO-AR       & 4.13 & 747  & 0.28 \\
seasonal naive & 5.62 & 1030 & 0.00 \\
climatology    & 8.59 & 1628 & $-0.58$ \\
\midrule
TSO (PSE)      & $\approx$2.0--2.6 & --- & --- \\
\bottomrule
\end{tabular}
\end{center}

\subsection{What each number says}

\begin{itemize}
  \item \textbf{Climatology $<$ naive.} Expected. The long-run average ignores
        the recent load level; last week carries far more signal than the
        all-time hourly mean. Losing to naive here is normal, not a bug.
  \item \textbf{Linear models remove $\sim$30\% of naive error.} Calendar +
        weather + weekly lags explain a lot. Ridge $\approx$ LASSO: with only
        $\sim$20 curated features there is little for LASSO's selection to prune.
        LASSO earns its keep when we throw hundreds of candidate features at it
        (the LEAR recipe for prices).
  \item \textbf{The 1.4 pp gap to the TSO is the project.} It will not close
        with linear models. Known missing pieces: nonlinear
        temperature--hour interactions (gradient boosting, M4), holiday-specific
        handling, and recency beyond day-old lags.
\end{itemize}

\subsection{Against the literature}

\begin{itemize}
  \item Grid operators typically report day-ahead MAPE below 3\%; competitive
        ISO practice sits near or below 2.5\%. PSE at $\approx$2--2.6\% on our
        own scoring is a normal, strong TSO.
  \item CAISO's day-ahead forecast ran 4.55\% MAPE in July 2024 while a
        commercial vendor scored 2.65\% on the same market --- TSO forecasts
        are beatable, but not by baselines.
  \item Academic national-level studies with tuned ML report 1--3\%.
        Our ridge at 4.0\% before any nonlinear model is on the expected path;
        gradient boosting typically lands in the 2.5--3.5\% range on
        country-level load with weather forecasts as input.
\end{itemize}

Sources: Hong \& Fan (2016), ``Probabilistic electric load forecasting:
A tutorial review'', \emph{Int.\ J.\ Forecasting} 32(3); CAISO/vendor
comparison reported by commercial forecasting benchmarks (2024); systematic
STLF review, \emph{Applied Energy} 326 (2022) 119851.

\subsection{Why MAPE is the headline (and what it hides)}

MAPE is scale-free and the industry's lingua franca --- PSE, ENTSO-E and
vendors all quote it. But it treats a 700~MW miss at 3~a.m.\ the same as at
the evening peak. RMSE (also in our table) punishes the big misses that cost
real balancing money. Rule: rank on skill-vs-naive, report MAPE for
comparability, watch RMSE and the per-hour breakdown for risk. Never MAPE
for prices (they cross zero; division explodes).

\subsection{Interview line}

``My baselines put the TSO 1.4 percentage points ahead of my best linear
model. That gap is exactly the value nonlinear models must prove --- and if
they close only half of it, the table says so.''
